\documentclass[aspectratio=169]{beamer}

\usepackage{xcolor}
\usepackage{listings}
\usepackage{booktabs}

\usetheme{Metropolis}

\title{Mathematics in slitex}
\subtitle{Equations, macros, and operators}
\author{Slitex}
\date{\today}

% Custom commands
\newcommand{\R}{\mathbb{R}}
\newcommand{\N}{\mathbb{N}}
\newcommand{\C}{\mathbb{C}}
\newcommand{\norm}[1]{\left\lVert #1 \right\rVert}
\newcommand{\abs}[1]{\left| #1 \right|}
\newcommand{\inner}[2]{\langle #1,\, #2 \rangle}

\DeclareMathOperator{\tr}{tr}
\DeclareMathOperator{\rank}{rank}
\DeclareMathOperator*{\argmax}{argmax}
\DeclareMathOperator*{\argmin}{argmin}

\begin{document}

\begin{frame}
  \titlepage
\end{frame}

% ── Inline math ───────────────────────────────────────────────────────────────

\begin{frame}{Inline Math}
  Use \texttt{\$...\$} for inline math:

  \begin{itemize}
    \item The energy $E = mc^2$ relates mass and energy.
    \item The golden ratio $\varphi = \frac{1 + \sqrt{5}}{2} \approx 1.618$.
    \item Euler's identity: $e^{i\pi} + 1 = 0$.
    \item For $x \in \R^n$, the norm $\norm{x} = \sqrt{x^\top x}$.
  \end{itemize}
\end{frame}

% ── Display equations ─────────────────────────────────────────────────────────

\begin{frame}{Display Equations}
  Numbered equation:
  \begin{equation}
    \mathcal{L}(\theta) = -\frac{1}{N} \sum_{i=1}^{N} \log p(y_i \mid x_i;\, \theta)
  \end{equation}

  Unnumbered equation:
  \begin{equation*}
    \hat{\theta} = \argmax_{\theta} \prod_{i=1}^{N} p(y_i \mid x_i;\, \theta)
  \end{equation*}
\end{frame}

% ── Aligned equations ─────────────────────────────────────────────────────────

\begin{frame}{Aligned Equations}
  \begin{align}
    (A + B)^\top  &= A^\top + B^\top       \\
    (AB)^\top     &= B^\top A^\top         \\
    (A^{-1})^\top &= (A^\top)^{-1}         \\
    \tr(AB)       &= \tr(BA)
  \end{align}
\end{frame}

% ── Custom macros ─────────────────────────────────────────────────────────────

\begin{frame}{Custom Macros via \texttt{\textbackslash newcommand}}
  \begin{itemize}
    \item \texttt{\textbackslash R} $\to$ $\R$ \quad (reals)
    \item \texttt{\textbackslash N} $\to$ $\N$ \quad (naturals)
    \item \texttt{\textbackslash C} $\to$ $\C$ \quad (complex)
    \item \texttt{\textbackslash norm\{x\}} $\to$ $\norm{x}$
    \item \texttt{\textbackslash abs\{z\}} $\to$ $\abs{z}$
    \item \texttt{\textbackslash inner\{u\}\{v\}} $\to$ $\inner{u}{v}$
  \end{itemize}

  \bigskip

  Macros with arguments use \texttt{\#1}, \texttt{\#2}, etc.\ inside the definition.
\end{frame}

% ── Math operators ────────────────────────────────────────────────────────────

\begin{frame}{Math Operators via \texttt{\textbackslash DeclareMathOperator}}
  Operators render in upright font and respect spacing:

  \begin{align*}
    \tr(A)    &= \sum_i A_{ii}                                              \\
    \rank(A)  &= \dim(\operatorname{im}(A))                                 \\
    w^*       &= \argmin_{w \in \R^d} \; \mathcal{L}(w)                    \\
    \hat{y}   &= \argmax_{c \in \mathcal{C}} \; p(c \mid x)
  \end{align*}

  \bigskip

  The \texttt{*} variant (\texttt{\textbackslash DeclareMathOperator*}) places limits
  below the operator in display mode.
\end{frame}

% ── Gather & multline ─────────────────────────────────────────────────────────

\begin{frame}{Gather \& Multiline}
  \texttt{gather} — centered, separate lines:
  \begin{gather*}
    f(x) = ax^2 + bx + c \\
    g(x) = \sin(x) + \cos(x) \\
    h(x) = e^{-x^2/2}
  \end{gather*}
\end{frame}

\end{document}
